\section{Impacto social: deflación del trabajo de cuello blanco y ventana de desempleo}

La otra mitad del tono de este episodio proviene del impacto social. Los modelos de Coding multiplican por diez, o incluso por varias decenas, la productividad de los mejores investigadores e ingenieros; los ciclos de investigación e ingeniería se comprimen y las organizaciones, de manera natural, recalculan sus necesidades de personal. La deflación del trabajo de cuello blanco y la ventana de desempleo no son ciencia ficción lejana, sino un problema real: la capacidad de absorción de las organizaciones no alcanza a seguir el salto de productividad.

\begin{figure}[H]
\centering
\includegraphics[width=0.92\textwidth]{social-impact-chain.png}
\caption{Cómo el acelerador de Coding se transmite a la deflación del trabajo de cuello blanco y a la presión sobre el empleo.}
\end{figure}

\begin{knowledgebox}{Lectura de la figura: la cadena del impacto social}
La cadena comienza con el salto de los modelos de Coding, se transmite a la compresión de los ciclos de investigación y desarrollo, luego al cambio en la demanda de las organizaciones y, por último, se convierte en presión sobre los salarios y los puestos de trabajo. Lo importante no es la frase vacía «¿la IA reemplazará a las personas?», sino si la velocidad de difusión de la tecnología supera la velocidad con la que las organizaciones y las instituciones sociales pueden absorberla.
\end{knowledgebox}

\subsection{Por qué el trabajo de cuello blanco lo resiente primero}

Muchas tareas de cuello blanco, como Coding, redacción, análisis, producto, operaciones e investigación, ocurren en el mundo digital: sus entradas y salidas pueden expresarse como texto y código, por lo que es más fácil conectarlas a los modelos. Las tareas manuales y del mundo físico están limitadas por la robótica, el hardware y la complejidad del entorno, así que su difusión puede ser más lenta.

\begin{warningbox}{El aumento de productividad no se traduce automáticamente en mayor bienestar}
Si las ganancias de productividad se concentran en unas pocas empresas, en unos cuantos individuos sobresalientes o en el lado del capital, los trabajadores de cuello blanco comunes podrían experimentar primero presión salarial y reducción de puestos. Cómo se reparten los beneficios de la tecnología es un problema social, no algo que la capacidad de los modelos resuelva por sí sola.
\end{warningbox}

\subsection{Nuevas reflexiones sobre inversión}

En el plano de la inversión, este episodio contiene un juicio implícito: si Coding es un acelerador, las empresas de modelos y de cadenas de herramientas se revalorizarán; si el modelo se convierte en el OS, también se reorganizarán los puntos de entrada de las aplicaciones y el panorama de plataformas. Sin embargo, invertir no puede basarse solo en palabras de moda: hay que ver si el modelo entra en flujos de trabajo reales, si genera retroalimentación sostenible y si tiene ventajas de costo y canales de distribución.

\subsection{Resumen de la sección}

El impacto social de Coding/Agent proviene de la compresión de la productividad. Traerá nuevas empresas, nuevos puntos de entrada y nuevas oportunidades de inversión, pero también una revaloración de los puestos de cuello blanco, presión salarial y desafíos institucionales.
